\section{The quantum entropy function and the logarithmic corrections}\label{sec:quantum}

\subsection{The quantum entropy function}\label{ssec:qef}

Sen \cite{Sen08qef} proposed that the exact degeneracy of the horizon states of an extremal black hole with charges $\vec q$ is the finite part of the Euclidean string theory path integral on the near-horizon $AdS_2$,
\begin{equation}\label{eq:qef}
d_{\rm hor}(\vec q) = \Big\langle \exp\Big[-i q_i\oint A^{(i)}_\theta\, d\theta\Big]\Big\rangle^{\rm finite}_{AdS_2} ,
\end{equation}
with a Wilson line insertion at the boundary carrying the charges, the boundary conditions fixing the charges rather than the chemical potentials, and the ``finite part'' defined by subtracting the terms proportional to the length of the boundary; and that this quantity equals the microscopic index $B_6$ for quarter-BPS black holes, after the hair modes are removed (Section~\ref{ssec:hair}).

\begin{verdict}[claim 27]
Confirmed. The proposal passes every test that can be computed: the classical limit of \eqref{eq:qef} is the exponential of the Wald entropy; the one-loop correction is the logarithmic term, which matches the microscopic index exactly in $\mathcal N=4$ and $\mathcal N=8$ (Theorem~\ref{thm:logs}); and the localisation of the path integral in $\mathcal N=2$ supergravity, proposed in \cite{BBGMS10} and carried out in \cite{DGM11,DGM13}, reproduces the exact Rademacher expansion of the index of $\mathcal N=8$ black holes and the leading Bessel function of the $\mathcal N=4$ one \cite{MurthyReys16,IMT25,IMT25b}. The last of these results is weaker than its statement in the literature, and Sen's own later assessment \cite{Sen23localisation} says so: \gap{the claim that supersymmetric localisation of the $AdS_2$ path integral reproduces the Rademacher expansion from first principles} rests on three steps that are not established. The Killing spinors of the near-horizon geometry grow at the boundary of $AdS_2$ and do not act on the space of normalisable fluctuations over which the path integral is defined, so the localisation argument has no proved domain; the measure on the localisation manifold used in \cite{DGM11,MurthyReys16} is postulated and not derived, and a different measure changes every subleading term; and the one-loop determinants over the transverse directions depend on an ultraviolet cut-off that the computation sets, without justification, at the four-dimensional Planck scale. Each of these is listed in \cite{Sen23localisation}, Sections 3 to 5 and 8, as open. There is no proof: the path integral is not defined beyond the semiclassical expansion, the localisation argument assumes an off-shell supersymmetric completion of the higher-derivative terms, and the identification of the subtracted terms with the hair is made case by case. The missing step is a definition of the right-hand side of \eqref{eq:qef} that does not rely on a saddle point expansion; the conjecture is, in this sense, a proposal for what that definition should give.
\end{verdict}

\subsection{The logarithmic corrections}\label{ssec:logs}

The one-loop contribution to \eqref{eq:qef} from the massless fields on $AdS_2\times S^2$ gives a term $c\ln A_H$ in the entropy, with $c$ determined by the spectrum and the two-derivative action alone. Banerjee, Gupta, Mandal and Sen \cite{BGMS11} computed it for quarter-BPS black holes in $\mathcal N=4$ and eighth-BPS black holes in $\mathcal N=8$ string theory, and Sen \cite{Sen12N2} for half-BPS black holes in $\mathcal N=2$ theories with $n_V$ vector and $n_H$ hypermultiplets.

\begin{theorem}[claim 28]\label{thm:logs}
The logarithmic correction to the entropy of a quarter-BPS black hole in any $\mathcal N=4$ string compactification is zero, and that of an eighth-BPS black hole in $\mathcal N=8$ string theory on $T^6$ is $-4\ln A_H$. Both agree with the microscopic indices: the asymptotic expansion of the Fourier coefficients of $1/\Phi_{10}$ and its CHL analogues has no logarithmic term, and that of the $\mathcal N=8$ index, a Fourier coefficient of $\vartheta_1(\tau,z)^2/\eta(\tau)^6$, has the term $-4\ln\Delta$ with $\Delta\propto A_H^2$, that is, $-2\ln\Delta = -4\ln A_H$ \cite{BGMS11}.
\end{theorem}

\begin{proof}
The macroscopic computation is the heat kernel of the quadratic fluctuation operator of all massless fields on $AdS_2\times S^2$ with the near-horizon flux, including the graviton, the gravitini and the vector fields, with the zero modes of the asymptotic symmetries treated separately; the coefficient is the $s^0$ term of the heat kernel together with the zero-mode counting, $-(c_{\rm local} + c_{\rm zm})\ln a$ with $a$ the radius. In $\mathcal N=4$ the contributions cancel between the gravity multiplet and the vector multiplets \cite[Section~7]{BGMS11}; in $\mathcal N=8$ the net result is $-8\ln a = -4\ln A_H$ \cite[Section~8]{BGMS11}. The microscopic expansions are in \cite[Section~2]{BGMS11}, from the saddle point expansion of the inverse Siegel form \cite{DavidSen06} and of the Jacobi form. The computation was rederived from the on-shell action of the Euclidean solution and from the $AdS_2$ Schwarzian in \cite{IMT25,IMT25b}, with the same coefficients.
\end{proof}

\begin{verdict}[claim 28]
Proved, as a match between two computations each of which is complete: the macroscopic coefficient is a theorem of one-loop quantum field theory on the near-horizon geometry and the microscopic one of the asymptotics of modular forms. The rotating black holes in four and five dimensions \cite{Sen12rotating} and the extremal Kerr and Reissner--Nordstr\"om black holes without supersymmetry \cite{Sen12N2,Sen13schwarzschild}, for which no microscopic count exists, are predictions and remain confirmed in the sense of claim~29.
\end{verdict}

\begin{proposition}[claim 29; Sen 2012]\label{prop:n2log}
For a half-BPS black hole in an $\mathcal N=2$ supergravity with $n_V$ vector and $n_H$ hypermultiplets, the logarithmic correction is $\frac{1}{12}(23 + n_H - n_V)\ln A_H$, equal to $\big(2 - \frac{\chi}{24}\big)\ln A_H$ with $\chi = 2(n_V - n_H + 1)$ the Euler number of the Calabi--Yau threefold in a type II compactification.
\end{proposition}

\begin{proof}
\cite{Sen12N2}, by the same heat kernel method, the only difference from $\mathcal N=4$ being the absence of the cancellation; the result was rederived by \cite{KLL14} from the heat kernel of $\mathcal N=2$ supergravity on the same background.
\end{proof}

\begin{verdict}[claim 29]
Confirmed. No microscopic count of $\mathcal N=2$ black hole states exists with which to compare; the comparison made in \cite{Sen12N2} is with the measure in the OSV formula, and agrees with the measure proposed by Denef and Moore if their result is extrapolated to weak topological string coupling. The missing step is a microscopic index for $\mathcal N=2$ black holes with an asymptotic expansion.
 The six papers that cite \cite{Sen12N2} and test a logarithmic correction do so in anti-de Sitter space or for the $\mathcal N=8$ theory; none supplies the microscopic count in $\mathcal N=2$ that the claim needs.
\end{verdict}
